\clearpage
\begingroup
\fancyhead[L]{\small REMARKS ON THE SEVENTH VOLUME.}
\renewcommand{\headrulewidth}{0pt}
\SourceLeaf{649}
\needspace{8\baselineskip}
\pdfbookmark[1]{Remarks on the Seventh Volume}{vii-volume-note}
\begin{center}\large REMARKS ON THE SEVENTH VOLUME.\end{center}
\begin{center}\rule{.09\linewidth}{.4pt}\end{center}

\par\noindent\hspace*{\parindent}The present seventh volume of \textsc{Gauss}’s Works opens with a reprint of the Theoria motus.
As is well known, \textsc{Schering} had already edited a new edition, published in 1871 by F. A. Perthes in a
format matching volumes I--VI of the collected edition, with some notes from the posthumous
papers added. Even so, it seemed advisable to include such a reprint in the
present seventh volume of the collected edition. A close examination made it possible to correct some
further inaccuracies in the first edition. The numerical examples were also
carefully recomputed, and a large number of errors newly found here were corrected
or, where that was impossible, at least identified. A large number of additional notes from
the posthumous papers relating to the work were also included.

\par\noindent\hspace*{\parindent}The volume also contains the entire previously unpublished collection of posthumous papers
by \textsc{Gauss} on theoretical astronomy, above all \textsc{Gauss}’s investigations of the perturbations of Pallas.
Specialists had eagerly awaited their publication for decades, and they will still be
of great interest today.

\par\noindent\hspace*{\parindent}Two shorter papers already included in Volume VI are reprinted
because their connection with the posthumous notes made it desirable to include them. These are the
papers on the zodiac of the celestial bodies, page 313, and on the table for parabolic motion,
page 371.

\par\noindent\hspace*{\parindent}Mr. \textsc{Brendel} in Göttingen was responsible for preparing and editing the volume.

\par\noindent\hspace*{\parindent}In preparing the perturbations of Ceres and Pallas, the editor had to supplement the
posthumous notes extensively. The notes for which generally only isolated formulas and numbers could
be taken from the posthumous papers were therefore set in petit type. To save
space, the larger numerical tables were also set in petit type, even where the numbers come from
the posthumous papers.

\par\noindent\hspace*{\parindent}Otherwise, as in the other volumes, the editor’s insertions are always in square
brackets [\,]. Notes from the correspondence and other sources that do not originate with \textsc{Gauss} are
enclosed in braces \(\bigl\{\bigr\}\).

\par\noindent\hspace*{\parindent}Further details about the printing of the individual notes can be found in the remarks
by the editor after each part.
\begin{flushright}\textsc{F. Klein}.\end{flushright}
\par\begin{center}\rule{.24\linewidth}{.4pt}\end{center}
\noindent\hfill81\textsuperscript{*}\par
\clearpage
\endgroup
